\subsection{Valuations, tautologies and consistency}
\label{sec:tautologies}
\slidesrc{Propositional logic: 19--21}

\subsubsection{Valuations: giving a truth value to a formula}
\slidesrc{Propositional logic: 19--20}

Given a formula, can you tell if it is true or not? \textbf{No}: you usually need to know
the truth values of its atomic propositions first. Is $p \lor q$ true? It depends on $p$
and $q$.

\begin{nbdef}{Valuation (slide 19)}
A \textbf{valuation} is a function $v$ that assigns a truth value (T or F) to each atomic
(primitive) proposition. In a programming language it would be a function
\code{Val(p: AtomicProp): BOOLEAN}. Given a valuation, we can compute for \emph{any}
formula whether it is true or false, by applying the truth tables bottom-up.
\end{nbdef}

\begin{nbexample}{evaluating $(p \lor q) \limp r$ (slide 20)}
Let $v(p) = \false$, $v(q) = \true$, $v(r) = \false$. Then
\begin{align*}
v\big((p \lor q) \limp r\big)
  &= (v(p) \lor v(q)) \limp v(r) && \reason{replace each letter by its value} \\
  &= (\false \lor \true) \limp \false && \\
  &= \true \limp \false && \reason{$\false \lor \true = \true$ (table of $\lor$)} \\
  &= \false && \reason{$\true \limp \false = \false$ (table of $\limp$)}
\end{align*}
So under this valuation the formula is \textbf{false}: the antecedent $p \lor q$ holds but
the consequent $r$ does not, which is the only way to falsify an implication.
\end{nbexample}

\begin{nbfact}{Rows = valuations}
Each row of a truth table corresponds to exactly one valuation of the atomic propositions
of the formula. A truth table therefore shows the value of the formula under
\emph{all} possible valuations ($2^n$ of them).
\end{nbfact}

\subsubsection{Tautologies, consistency, inconsistency}
\slidesrc{Propositional logic: 21}

Some formulas have interesting properties that do not depend on the particular valuation.

\begin{nbdef}{Tautology, consistent, inconsistent (slide 21)}
\begin{enumerate}[nosep]
  \item A formula is a \textbf{tautology} iff it is true under \emph{every} valuation.
  \item A formula is \textbf{consistent} iff it is true under \emph{at least one}
        valuation (also called \emph{satisfiable}).
  \item A formula is \textbf{inconsistent} iff it is \emph{not} made true under any
        valuation (also called \emph{contradiction} or \emph{unsatisfiable}).
\end{enumerate}
Since each row of a truth table is a valuation, truth tables decide these properties: look
at the last column. All T $\Rightarrow$ tautology; at least one T $\Rightarrow$ consistent;
no T $\Rightarrow$ inconsistent.
\end{nbdef}

\begin{center}
\begin{tikzpicture}[font=\small]
  \draw[thick, nbgreen, fill=nbgreen!8] (0,0) ellipse (4.2cm and 1.5cm);
  \draw[thick, nbblue, fill=nbblue!10] (1.5,0) ellipse (1.9cm and 0.95cm);
  \node[nbgreen!70!black, align=center] at (-1.9,0.15) {\textbf{consistent}\\(true in $\ge 1$ row)\\$p \lor q$,\ \ $p \land \neg q$};
  \node[nbblue, align=center] at (1.5,0) {\textbf{tautologies}\\(true in all rows)\\$p \lor \neg p$};
  \draw[thick, nbred, fill=nbred!8] (6.4,0) ellipse (1.7cm and 0.95cm);
  \node[nbred, align=center] at (6.4,0) {\textbf{inconsistent}\\(true in no row)\\$p \land \neg p$};
\end{tikzpicture}
\captionof{figure}{Every formula is either consistent or inconsistent. Tautologies are a
special case of consistent formulas. The negation of a tautology is inconsistent and vice
versa.}
\end{center}

\begin{nbexample}{one formula of each kind}
\begin{center}
\renewcommand{\arraystretch}{1.1}
\begin{tabular}{c|c|c|c}
$p$ & $p \lor \neg p$ & $p \land \neg p$ & $p \limp \neg p$ \\
\hline
F & T & F & T \\
T & T & F & F \\
\hline
 & \textcolor{nbblue}{tautology} & \textcolor{nbred}{inconsistent} & \textcolor{nbgreen!70!black}{consistent, not a tautology}
\end{tabular}
\end{center}
$p \lor \neg p$ (``it rains or it does not rain'') is always true: it carries no
information about the world. $p \land \neg p$ can never be true.
\end{nbexample}

\begin{nbexample}{modus ponens is a tautology, affirming the consequent is not}
\notinslides
Consider the two formulas $A = (p \land (p \limp q)) \limp q$ and
$B = (q \land (p \limp q)) \limp p$. Full truth table, one column per sub-formula:
\begin{center}
\renewcommand{\arraystretch}{1.1}
\begin{tabular}{cc|c|cc|cc}
$p$ & $q$ & $p \limp q$ & $p \land (p\limp q)$ & $A$ & $q \land (p \limp q)$ & $B$ \\
\hline
F & F & T & F & T & F & T \\
F & T & T & F & T & T & \cellcolor{nbred!20}F \\
T & F & F & F & T & F & T \\
T & T & T & T & T & T & T \\
\end{tabular}
\end{center}
$A$ is a \textbf{tautology}: whenever $p$ and $p \limp q$ are both true, $q$ is true. This
is the reasoning pattern \emph{modus ponens} (Section~\ref{sec:natural-deduction-prop}).
$B$ is \textbf{not}: in row 2 ($p$ false, $q$ true) both $q$ and $p \limp q$ are true but
$p$ is false. Concretely: ``if it rains, the street is wet; the street is wet; so it
rains'' is wrong, because someone may have washed the street.
\end{nbexample}

\subsubsection{Logical equivalence and logical consequence}
\notinslides

Two more notions, which the slides use implicitly (slide 16 and the proof-theory slides),
are worth stating precisely.

\begin{nbdef}{Logical equivalence}
Two formulas $\varphi$ and $\psi$ are \textbf{logically equivalent}, written
$\varphi \equiv \psi$, iff they have the same truth value under \emph{every} valuation,
i.e.\ iff $\varphi \liff \psi$ is a \textbf{tautology}.
\end{nbdef}

\begin{nbwarning}[About slide 16]
Slide 16 says: ``If $p \liff q$ is true, then $p$ and $q$ are said to be logically
equivalent''. Read it carefully: what is meant is that $p \liff q$ is true under
\emph{every} valuation (a tautology). If $p \liff q$ is true under \emph{one} valuation,
$p$ and $q$ merely happen to have the same value there. Example: $p$ = ``it is Monday'' and
$q$ = ``it rains'' may both be true today, so $p \liff q$ is true today, but they are
certainly not logically equivalent.
\end{nbwarning}

\begin{nbdef}{Logical consequence (valid argument)}
A formula $\psi$ is a \textbf{logical consequence} of the premises
$\varphi_1, \ldots, \varphi_n$, written $\varphi_1, \ldots, \varphi_n \entails \psi$, iff
\emph{every valuation that makes all the premises true also makes $\psi$ true}. Equivalently:
there is no valuation with all premises true and $\psi$ false; equivalently,
$(\varphi_1 \land \cdots \land \varphi_n) \limp \psi$ is a tautology. An argument whose
conclusion is a logical consequence of its premises is called \textbf{(deductively)
valid}.
\end{nbdef}

\begin{nbexample}{a few useful equivalences, checked by truth table}
\begin{center}
\renewcommand{\arraystretch}{1.1}
\begin{tabular}{cc|cc|cc|cc}
$p$ & $q$ & $p \limp q$ & $\neg p \lor q$ & $\neg q \limp \neg p$ & & $\neg(p \land q)$ & $\neg p \lor \neg q$ \\
\hline
F & F & T & T & T & & T & T \\
F & T & T & T & T & & T & T \\
T & F & F & F & F & & T & T \\
T & T & T & T & T & & F & F \\
\end{tabular}
\end{center}
Equal columns $\Rightarrow$ equivalent formulas:
$p \limp q \equiv \neg p \lor q \equiv \neg q \limp \neg p$ (the last one is the
\emph{contrapositive}), and $\neg(p \land q) \equiv \neg p \lor \neg q$ (one of De Morgan's
laws, used again in Section~\ref{sec:natural-deduction-prop}).
\end{nbexample}

\begin{nbwhy}[Truth tables in real software]
Deciding whether a formula is consistent (``is there a valuation that makes it true?'') is
the famous \textbf{SAT problem}. A truth table needs $2^n$ rows, which explodes quickly
(the combinatorial explosion of Section~\ref{sec:ai-history}); modern \textbf{SAT solvers}
use clever search instead and routinely handle formulas with millions of variables. They
are used to verify hardware circuits, find bugs in software, schedule exams and timetables,
and solve puzzles such as Sudoku (``cell (1,1) contains 5'' is an atomic proposition).
\end{nbwhy}
